\documentclass[11pt,a4paper]{article}
\usepackage[T1]{fontenc}
\usepackage[utf8]{inputenc}
\usepackage[english]{babel}
\usepackage{amsmath,amssymb}
\usepackage[margin=2cm]{geometry}
\setlength{\parindent}{0pt}
\setlength{\emergencystretch}{3em}
\allowdisplaybreaks

\title{About \texttt{book\_irrational\_equations.tex}}
\author{}
\date{}
\begin{document}
\maketitle
This note explains what the test data in \texttt{book\_irrational\_equations.tex} is used for. That file is input for the engine in \texttt{main.c}; it is read by the test \texttt{test\_irrational\_equations.c} in this folder.

\begin{enumerate}
\item \textbf{Purpose.} The file holds 10 true statements. Each says that an equation with the unknown under one or more square roots holds exactly for the values printed as its result in a school textbook, for example $\sqrt{x} - 3 = \sqrt{3} - \sqrt{12 - x} \iff x = 3 \lor x = 9$.

\item \textbf{What it exercises.} The sign table with square roots (rule 18p-7): each root is a letter $s$ with $s^{2}$ equal to its radicand, the sign of $P + Qs$ comes from the signs of $P$, $Q$ and $P^{2} - Q^{2}R$, and a comparison has a value only where every radicand is not negative. Two and three roots in the same equation, and roots that are multiples of one another.

\item \textbf{Variety.} One, two and three square roots, radicands of first and second degree, a root in a denominator, equations whose squared form has a root that does not satisfy the original one (it must not appear among the solutions).

\item \textbf{Expected result.} The line \texttt{PASS: 10 of 10 statements proved, each with a written proof}.

\item \textbf{How to run it.} From the top folder of the repository:
\begin{verbatim}
cc -std=c11 -O2 -Wall -Wextra -o /tmp/test_irrational_equations \
    tests/equations_inequalities/test_irrational_equations.c \
    && /tmp/test_irrational_equations
\end{verbatim}
\end{enumerate}
\end{document}
